% Figure: free decay under dry (Coulomb) friction versus viscous damping.
% Coulomb decay has a straight-line (arithmetic) envelope; viscous decay
% has an exponential (geometric) envelope. Both start at the same peak.
% The two envelopes are drawn as guide lines and one oscillating trace
% is shown for the Coulomb case.
\begin{figure}[htbp]
\centering
\begin{tikzpicture}
\begin{axis}[
  width=12cm, height=7cm,
  xlabel={time $t$ (periods)}, ylabel={displacement $x$},
  xmin=0, xmax=8, ymin=-1.15, ymax=1.15,
  axis lines=left,
  domain=0:8, samples=600,
  legend pos=north east,
  legend cell align=left,
  every axis x label/.style={at={(current axis.right of origin)}, anchor=west},
  every axis y label/.style={at={(current axis.above origin)}, anchor=south},
  grid=major, grid style={enggray!20},
  tick label style={font=\scriptsize},
  label style={font=\small},
  legend style={font=\scriptsize},
]
% Coulomb: envelope falls by a fixed step 4 F_f / k each cycle. Take the
% per-cycle drop as 0.125 (so amplitude reaches zero near t = 8).
% Oscillation at unit period, envelope (1 - 0.125 t) clamped at zero.
\addplot[engblue, very thick, samples=800]
  {max(0, 1 - 0.125*x) * cos(deg(2*pi*x))};
\addlegendentry{Coulomb (dry friction)}
% Coulomb envelope, straight lines
\addplot[engblue, thin, dashed, forget plot] {max(0, 1 - 0.125*x)};
\addplot[engblue, thin, dashed, forget plot] {-max(0, 1 - 0.125*x)};
% Viscous envelope for comparison, same initial amplitude, exponential
\addplot[engred, thick, dashed] {exp(-0.30*x)};
\addlegendentry{viscous envelope}
\addplot[engred, thick, dashed, forget plot] {-exp(-0.30*x)};
% dead-band note
\node[engblue, font=\scriptsize, anchor=west] at (axis cs:6.15,0.28)
  {reaches rest, not zero};
\end{axis}
\end{tikzpicture}
\caption[Dry-friction free decay against viscous decay]{Free decay of an
oscillator damped by dry (Coulomb) friction, blue, compared with the
exponential envelope of viscous damping, red, both released from the
same peak. The Coulomb envelope falls by a fixed increment $4F_f/k$
every cycle, so its bound is a straight line that reaches zero in a
finite number of cycles, while the viscous envelope decays by a fixed
ratio each cycle and only approaches zero asymptotically. A dry-friction
oscillator halts wherever the spring force drops below the static
friction band rather than settling at the geometric origin, the
stick-band that leaves a residual offset the viscous model never shows.
The frequency of a Coulomb-damped oscillator stays at the undamped
$\omega_n$ throughout, since dry friction does not shift the natural
frequency the way viscous damping shifts it to $\omega_d$.}
\label{fig:vol03ch06:coulomb-decay}
\end{figure}
